\section{Conclusion}

\texttt{voiage} treats Value of Information analysis as a comparison of
avoidable decision losses rather than a single calculation. The same decision
may be affected by uncertain parameters, the proposed evidence, model
structure, population differences, implementation, or timing. Keeping these
sources distinct helps show whether the relevant response is further research,
better delivery, a different model, or a different decision.

Version 1.0.0 provides established EVPI, EVPPI, EVSI, ENBS, dominance, CEAF,
heterogeneity, and structural analyses, alongside newer methods whose status
and validation limits remain visible. The worked health example illustrates
how evidence priorities and study value depend on the declared sampling,
population, cost, delay, and uptake assumptions. Applied use still depends on
a defensible decision model and evidence relevant to the target setting.
